\section{Opposite multiplicity excesses in every prescribed phase box}
\label{sec:localized}

Fix $q\ge1$ and an odd prime $d\nmid q$. Put $h=\varphi(q)$ and
$\kappa=(\cdot/d)$. For each character $\chi\bmod q$ use its primitive
inducing character $\chi^*$, of conductor $f_\chi$. The character
$\chi^*\kappa$ is primitive of conductor $df_\chi$. At positive height put
\[
 z(\gamma)=\sum_{\chi\bmod q}
 \left(\operatorname{ord}_{1/2+i\gamma}L(s,\chi^*)
      -\operatorname{ord}_{1/2+i\gamma}L(s,\chi^*\kappa)\right),
 \qquad z_\pm=\max(\pm z,0).
\]
Omitted Euler factors have their zeros on $\Re s=0$. Thus the same integer
is obtained with the imprimitive functions. In the range of
\texttt{isPolar\_primeRemainder} this is precisely its residue numerator.
Write $\mathbb T=\mathbb R/(2\pi\mathbb Z)$. For a finite set $B$ of
primes put $\theta_B(t)=(t\log b\bmod 2\pi)_{b\in B}$.

\paragraph{Theorem.}
Let $U$ be a nonempty open box of $\mathbb T^B$, and fix $0<\alpha\le1$.
With Dirichlet GRH as established in the source, there are positive
constants $c,T_0$ depending only on $q,d,B,U,\alpha$ such that, simultaneously
for both signs, $T\ge T_0$ and $T^\alpha\le H\le T$,
\begin{equation}
 \sum_{\substack{T<\gamma\le T+H\\\theta_B(\gamma)\in U}}z_\pm(\gamma)
    \ \ge\ c H\log\log T .
 \label{eq:phase-excess}
\end{equation}
Each distinct ordinate is summed once. The left side counts excess
multiplicity after all coincidences, including coincidences across
different characters, have been cancelled. The case $B=\varnothing$
gives the unrestricted statement. For $q=1$ these are respectively the
zeros and poles, with multiplicities, of $\zeta(s)/L(s,\kappa)$.

\paragraph{A uniform rational-frequency estimate.}
For a primitive character $\xi$ of fixed conductor (including the
zeta case), put
$D_\xi(x)=\sum_{T<\gamma_\xi\le T+H}x^{i\gamma_\xi}$.
If $x=a/b>1$ is in lowest terms with $a,b\le Q$, $Q\ge2$, then
\begin{equation}
 D_\xi(x)=-\frac{H}{2\pi}\frac{\Lambda(x)\xi(x)}{\sqrt{x}}
              +O_\xi(Q^4\log^2(2T)).
 \label{eq:rational-landau}
\end{equation}
Here the numerator in the main term is zero unless $x$ is an integral
prime power. The value at $x^{-1}$ is the conjugate. Banks proves a
sharper uniform estimate for $x\ge2$ in Theorem 1.2 of
\emph{Comparing zeros of distinct Dirichlet $L$-functions}
(\href{https://arxiv.org/abs/2302.07073}{arXiv:2302.07073}).
We give the coarser estimate for the full range $x>1$ because some
ratios of products of primes approach~$1$.

Integrate $x^s L'(s,\xi)/L(s,\xi)$, divided by $2\pi i$, on the
rectangle with vertical sides $-1,2$. Move each horizontal side by
at most $1$ to a height at distance $\gg_\xi 1/\log(2T)$ from all
nearby zero ordinates. The standard local zero count
$N_\xi(t+1)-N_\xi(t-1)\ll_\xi\log(2t)$ permits this choice.
The local partial-fraction expansion for $L'/L$ then gives
$O_\xi(\log^2(2T))$ on the two horizontal sides (Banks, Lemma 2.2).
The horizontal integrals cost $O_\xi(Q^2\log^2(2T))$. Restoring the
endpoints costs $O_\xi(\sqrt{x}\log(2T))$ in the sum of $x^\rho$ by
GRH. On the right side expand $L'/L=-\sum_{n\ge2}\Lambda(n)\xi(n)n^{-s}$.
The diagonal $n=x$ gives the displayed main term before division by
$\sqrt{x}$. For $n\ne x$ the oscillatory integral of $(x/n)^{it}$ is
$O(1/|\log(x/n)|)$. For $x/2<n<2x$ we have $|x-n|\ge1/Q$ and
$|\log(x/n)|\gg Q^{-2}$. Elsewhere the reciprocal is bounded. Since
$\sum_{n\ge2}\Lambda(n)/n^2<\infty$, the total off-diagonal cost is
$O(Q^4)$. On the left side the functional equation and Stirling's
formula give a logarithmic term $-\log(f_\xi t/(2\pi))$, a uniformly
$O_\xi(1/t)$ remainder, and a logarithmic derivative at $2-it$.
Integration by parts bounds the logarithmic term by
$O_\xi(\log(2T)/\log x)$, with $1/\log x\ll Q$. The absolutely
convergent Dirichlet series at $2-it$ has frequencies $\log(nx)\ge
\log 2$, so costs $O(1)$. These estimates prove
\eqref{eq:rational-landau} after division by $\sqrt{x}$. At frequency
$1$ use the Riemann--von Mangoldt formula. Subtracting the equal numbers
of degree-one functions cancels the common $t\log t$ term and gives
\begin{equation}
  \sum_{T<\gamma\le T+H} z(\gamma)
    =-\frac{hH}{2\pi}\log d+O_{q,d}(\log(2T)).
 \label{eq:conductordifference}
\end{equation}

\paragraph{Choosing a bounded test with a growing prime contribution.}
Let $Y=c_0\log T$ for a sufficiently small fixed $c_0>0$ to be specified
by the error bound below. Let $S$ consist of primes $p\le Y$ outside
$B$ satisfying $p\equiv1\pmod q$ and $\kappa(p)=-1$.
Define $w_p=(\log p)/\sqrt p$, $\sigma^2=\sum_{p\in S}w_p^2$.
This union of reduced classes modulo $qd$ has density $1/(2h)$ among
primes. The prime number theorem in fixed progressions and partial
summation give
\[
 \sigma^2\sim\frac{(\log Y)^2}{4h}.
\]
For large $T$, define the real trigonometric polynomial
\[
 A(t)=\Im\prod_{p\in S}
       \left(1+i\frac{w_p}{\sigma}\cos(t\log p)\right).
\]
Pointwise $|A|\le\exp(\sum_p w_p^2/(2\sigma^2))=\sqrt e$.
Its constant term is zero. The coefficient of
$\cos(t\log p)$ arising from the singleton $p$ is $w_p/\sigma$.
All remaining nonzero terms involve at least two distinct primes
(in fact at least three), each with exponent $1$ or $-1$ in the
frequency ratio. Hence neither this ratio nor its inverse is an integral
prime power. At a singleton $p$ we have
$\sum_\chi\chi^*(p)(1-\kappa(p))=2h$.
Taking the real part in \eqref{eq:rational-landau} consequently gives the
weighted singleton contribution $-hH\sigma/\pi$.

\paragraph{Restricting the phases.}
Take a smooth function $0\le F\le1$ supported inside $U$ with normalized
Haar integral $m>0$ (for empty $B$ take $F=1$). Write $F_M$ for its
rectangular Fourier truncation to exponents $|k_b|\le M$, where
$M=\lfloor\sqrt{\log T}\rfloor$. Smoothness gives
$\|F-F_M\|_\infty\ll_{F,D}M^{-D}$ for every fixed $D$ and a uniformly
bounded sum of absolute Fourier coefficients. The sum of $|z|$ in the
window is $O_{q,d}(H\log T+\log T)$ by the zero count. Thus $F_M$ can
replace $F$ in sums against $z$ and $zA$ with an error $o(H)$.
Here and below $F$ in a sum over ordinates stands for $F(\theta_B(\gamma))$.

In $F_M A$ every frequency with a nonzero exponent in $B$ also has
a nonzero exponent in $S$. Unique factorization makes each such
frequency, and its reciprocal, ineligible for a prime-power main term.
Only the constant Fourier coefficient $m$ of $F_M$ couples to a
singleton in $A$. For $F_M$ alone the constant frequency contributes
$O_{q,d}(H+\log T)$ by \eqref{eq:conductordifference}. Its nonconstant
prime-power terms contribute $O_{F,q,d,B}(H)$, since the factor
$\log b/b^{|k|/2}$ is bounded for the fixed primes $b\in B$ and the
Fourier coefficients are absolutely summable.

For completeness, all the rational frequencies occurring here, in
lowest terms, have numerator and denominator at most
\[
 Q=\exp\!\left(\sum_{p\le Y}\log p+M\sum_{b\in B}\log b\right)
       \le \exp(2Y+O_B(M))
\]
for large $T$. The sum of absolute Fourier coefficients of $A$ is at
most $\prod_{p\in S}(1+w_p/\sigma)\le\exp(\sqrt{|S|})$ by Cauchy--Schwarz.
Using \eqref{eq:rational-landau} termwise therefore costs
$O_{F,q,d}(Q^4\log^2(2T)\exp(\sqrt{|S|}))$. Choose, for example,
$c_0=\alpha/32$. This error is at most $T^{\alpha/4+o(1)}=o(H)$.
We have proved the two estimates
\[
 \sum_{T<\gamma\le T+H}z(\gamma)F=O(H),\qquad
 \sum_{T<\gamma\le T+H}z(\gamma)F A(\gamma)
        =-\frac{mh}{\pi}H\sigma+o(H).
\]
The functions $F(\sqrt e-A)$ and $F(\sqrt e+A)$ take their values
in $[0,2\sqrt e]$ and vanish outside the box. Their signed sums are,
respectively, $+(mh/\pi)H\sigma+O(H)$ and
$-(mh/\pi)H\sigma+O(H)$. Dividing the first by $2\sqrt e$ bounds the
positive mass from below, and negating the second does the same for the
negative mass. Since $\sigma\asymp\log\log T$,
\eqref{eq:phase-excess} follows.

\paragraph{Relation to earlier work.}
The signed list $z(\gamma)/h$ already occurs in the ProveIt archive
\texttt{grh-torsion-and-zero-data.zip}, where it is used to recover
$q$ and $d$. The present question asks where its two signs persist.
Liu and Zaharescu study smooth tests on multidimensional tori in races
between the zeros of two Dirichlet functions
(\href{https://doi.org/10.1016/j.jmaa.2020.124591}{JMAA 494 (2021), 124591}).
Banks compares zero multisets in short height intervals.
Bombieri and Perelli's \emph{Distinct zeros of $L$-functions}
(\href{https://doi.org/10.4064/aa-83-3-271-281}{Acta Arith. 83 (1998), 271--281})
treats large global counts for pairs of $L$-functions under specified
hypotheses. Here the cancellation is between two whole products,
the angular box is fixed arbitrarily before the amplifier is chosen,
and the same height window works in both directions.
This proof makes no priority claim. Its count allows different
multiplicities at a common zero. Simplicity, disjointness of zero sets,
and the conjectural $H\log T$ scale for almost all excess multiplicity
are stronger assertions.
